% Lösungen Brüche und Dezimalzahlen (zusammenbau v0.4, Heft)
\einheitenkopf{Brüche und Dezimalzahlen}

\erg{4}{a) $\frac{3}{9}$ \quad b) $\frac{3}{8}$ \quad c) $\frac{4}{12} = \frac{1}{3}$ \quad d) $\frac{4}{16} = \frac{1}{4}$ \quad e) $\frac{13}{25} = \frac{52}{100} = 52\,\%$ \quad f) $\frac{11}{20} = \frac{55}{100} = 55\,\%$ \quad g) $6$ Kästchen \quad h) $10$ Kästchen \quad i) $8$ Kästchen \quad j) $9$ Kästchen \quad k) $20 : 5 \cdot 4 = 16$ Kästchen \quad l) $18 : 3 \cdot 2 = 12$ Kästchen \quad m) $5$ Kästchen \quad n) $4$ Kästchen \quad o) Beide Diagonalen und die beiden Verbindungen gegenüberliegender Seitenmitten teilen das Quadrat in acht gleiche Dreiecke; eines davon kennzeichnen. \quad p) Zwei zueinander senkrechte Durchmesser teilen den Kreis in vier gleiche Viertel; drei davon kennzeichnen. \quad q) Q, denn $\frac{60}{360} = \frac{1}{6}$ \quad r) S, denn $\frac{72}{360} = \frac{1}{5}$}

4g)

\bruchrechteck{6}{20}


4h)

\bruchrechteck{10}{25}


4i)

\bruchrechteck{8}{32}


4j)

\bruchrechteck{9}{36}


4k)

\bruchrechteck{16}{20}


4l)

\bruchrechteck{12}{18}


4m)

\bruchrechteck{5}{12}


4n)

\bruchrechteck{4}{9}


4p)

\kreissektor{270}{}

\erg{5}{a) $120 : 4 \cdot 3 = 90$ l sind drin, es fehlen $120 - 90 = 30$ l \quad b) $1\,500 : 5 \cdot 2 = 600$ l sind drin, es passen noch $1\,500 - 600 = 900$ l hinein \quad c) $2{,}5 : 5 \cdot 3 = 1{,}5$ kg \quad d) $1{,}6$ kg $= 1\,600$ g, $1\,600 : 8 \cdot 3 = 600$ g}
\erg{6}{a) $\frac{12}{16} = \frac{3}{4} = 75\,\%$ \quad b) $\frac{10}{25} = \frac{2}{5} = 40\,\%$}
\erg{7}{a) z. B. $0{,}5$, denn $\frac{1}{4} = 0{,}25$ und $\frac{3}{5} = 0{,}6$; jede Zahl dazwischen ist richtig \quad b) z. B. $0{,}6$, denn $\frac{2}{5} = 0{,}4$ und $\frac{3}{4} = 0{,}75$; jede Zahl dazwischen ist richtig}
\erg{8}{a) $\frac{13}{5} > \frac{5}{2}$, denn $\frac{13}{5} = 2{,}6$ und $\frac{5}{2} = 2{,}5$; $\sqrt{5} \approx 2{,}24$ \quad b) $\frac{9}{5} > \frac{7}{4}$, denn $\frac{9}{5} = 1{,}8$ und $\frac{7}{4} = 1{,}75$; $\sqrt{3} \approx 1{,}73$ \quad c) ausgelassen \quad d) ausgelassen}
\erg{9}{a) $0{,}5^2 = 0{,}25$; die anderen sind $5{,}5$ und zweimal $0{,}55$ \quad b) $45\,\% = 0{,}45$; die anderen sind $0{,}7^2 = 0{,}49$, $0{,}5$ und $4{,}9$ \quad c) $-\frac{3}{4}$; $-0{,}749$; $1{,}7$; $\sqrt{3}$, denn $-\frac{3}{4} = -0{,}75$ und $\sqrt{3} \approx 1{,}73$ \quad d) $-0{,}41$; $-\frac{2}{5}$; $2{,}2$; $\sqrt{5}$, denn $-\frac{2}{5} = -0{,}4$ und $\sqrt{5} \approx 2{,}24$ \quad e) $\frac{12}{5} > \frac{9}{4}$, denn $\frac{12}{5} = 2{,}4$ und $\frac{9}{4} = 2{,}25$; $\sqrt{5} \approx 2{,}236$ \quad f) $\frac{17}{5} > \frac{13}{4}$, denn $\frac{17}{5} = 3{,}4$ und $\frac{13}{4} = 3{,}25$; $\sqrt{10} \approx 3{,}16$ \quad g) ausgelassen \quad h) ausgelassen \quad i) $(-0{,}8 + (-0{,}7)) : 2 = -0{,}75$ \quad j) $(-1{,}3 + (-1{,}2)) : 2 = -1{,}25$}
